\chapter{Unabhängigkeit unter verschiedenen Semantiken}
\label{ch:semantic-independence}

Dieses Kapitel untersucht bedingte Unabhängigkeit unter der Complete-,
Stable- und Preferred-Semantik in einer gemeinsamen Darstellung. Zunächst
wird der abstrakte Unabhängigkeitsbegriff auf die Labellings eines
Argumentationssystems übertragen. Anschließend werden einheitliche
aussagenlogische Bausteine entwickelt, aus denen sich die drei
Semantikkodierungen und eine gemeinsame QBF-Kodierung der Unabhängigkeit
zusammensetzen. Die Quantorenstruktur dieser Formeln erklärt zugleich die
unterschiedliche Komplexität der drei Entscheidungsprobleme. Abschließend
wird das Vorkommen von Unabhängigkeit in einer Auswahl von AFs empirisch untersucht.

Die semantische Definition geht auf
\textcite{rienstra:2020:independence-af} zurück. Notation und
Komplexitätsergebnisse orientieren sich an
\textcite{bluemel:2025:independence-aba}; QBFs als Zielformalismus für
Argumentationsprobleme wurden insbesondere von
\textcite{egly:woltran:2006:reasoning-qbf} untersucht.

\section{Semantikspezifische Unabhängigkeit}
\label{sec:sigma-independence-af}

Die abstrakte Unabhängigkeitsrelation wird im Folgenden durch die unter einer
Semantik zulässigen Labellings bestimmt.

\begin{Definition}
    \label{def:sigma-independence-af}
    Sei \(F=(A,R)\) ein Argumentationssystem, sei \(\sigma\) eine Semantik
    und seien \(X,Y,Z\subseteq A\) paarweise disjunkt. Die Menge \(X\) ist
    \emph{\(\sigma\)-unabhängig von \(Y\) gegeben \(Z\) in \(F\)},
    geschrieben
    \[
        X\perp_\sigma Y\mid Z,
    \]
    genau dann, wenn für alle
    \(\lambda_1,\lambda_2\in\Lambda_\sigma(F)\) aus
    \(\lambda_1|_Z=\lambda_2|_Z\) die Existenz eines
    \(\lambda_3\in\Lambda_\sigma(F)\) folgt, für das
    \[
        \lambda_3|_X=\lambda_1|_X,
        \qquad
        \lambda_3|_Y=\lambda_2|_Y,
        \qquad
        \lambda_3|_Z=\lambda_1|_Z=\lambda_2|_Z
    \]
    gilt. Andernfalls sind \(X\) und \(Y\) \emph{\(\sigma\)-abhängig
    gegeben \(Z\) in \(F\)}; hierfür schreiben wir
    \(X\not\perp_\sigma Y\mid Z\).

    Für \(\sigma\in\{\mathsf{co},\mathsf{st},\mathsf{pr}\}\) bezeichnet
    \(\mathrm{IND}_\sigma\) das folgende \emph{Entscheidungsproblem}:
    \[
        \begin{array}{ll}
        \text{Eingabe:}
            &F=(A,R)\text{ sowie paarweise disjunkte Mengen }
             X,Y,Z\subseteq A,\\
        \text{Frage:}
            &\text{Gilt }X\perp_\sigma Y\mid Z\text{ in }F\,?
        \end{array}
    \]
    Für \(x,y\in A\) stehen
    \(x\perp_\sigma y\mid Z\) und
    \(x\not\perp_\sigma y\mid Z\) kurz für die entsprechenden Aussagen
    über die Singletonmengen \(\{x\}\) und \(\{y\}\).
\end{Definition}

Die Definition verlangt nicht nur, dass die gewählten Bewertungen von \(X\)
und \(Y\) jeweils vorkommen. Sie verlangt ihre gemeinsame Realisierbarkeit
unter derselben Bewertung der Bedingungsmenge \(Z\).

\begin{Lemma}
    \label{lem:independence-trivial-labelling-set}
    Gilt \(|\Lambda_\sigma(F)|\leq 1\), so gilt
    \(X\perp_\sigma Y\mid Z\) für alle paarweise disjunkten Mengen
    \(X,Y,Z\subseteq A\).
\end{Lemma}

\begin{proof}
    Für \(\Lambda_\sigma(F)=\emptyset\) ist die Allquantifizierung über
    \(\lambda_1\) und \(\lambda_2\) leer. Besteht
    \(\Lambda_\sigma(F)\) aus genau einem Labelling \(\lambda\), so kann
    für das einzige Paar \((\lambda,\lambda)\) dieses Labelling selbst als
    \(\lambda_3\) gewählt werden.
\end{proof}

Das Lemma ist für die spätere empirische Auswertung wesentlich: Eine hohe
Unabhängigkeitsrate kann sowohl durch echte Kombinierbarkeit als auch durch
das Fehlen mehrerer vergleichbarer Labellings entstehen.

\section{Einheitliche Kodierung der Semantiken}
\label{sec:uniform-semantics-encoding}

Für alle drei Semantiken wird dieselbe dreiwertige Darstellung verwendet.
Dies vermeidet semantikspezifische Gleichheitsformeln und erlaubt eine
gemeinsame Unabhängigkeitskodierung.

\subsection{Labellingvariablen und Wohlgeformtheit}

\begin{Definition}
    \label{def:labelling-variable-block}
    Sei \(F=(A,R)\) ein Argumentationssystem. Ein
    \emph{Labellingvariablenblock} für \(F\) ist ein Tripel
    \(\mathbf V=(I,O,U)\) paarweise disjunkter Variablenmengen
    \[
        I=\{i_a\mid a\in A\},
        \qquad
        O=\{o_a\mid a\in A\},
        \qquad
        U=\{u_a\mid a\in A\}.
    \]
    Die Variablen \(i_a\), \(o_a\) und \(u_a\) kodieren die Labels
    \(\mathsf{in}\), \(\mathsf{out}\) beziehungsweise \(\mathsf{und}\)
    von \(a\). Die Formel
    \[
        \operatorname{One}_F(\mathbf V)
        :=
        \bigwedge_{a\in A}
        \Bigl(
            (i_a\lor o_a\lor u_a)
            \land\neg(i_a\land o_a)
            \land\neg(i_a\land u_a)
            \land\neg(o_a\land u_a)
        \Bigr)
    \]
    heißt \emph{Wohlgeformtheitsformel} des Blocks.

    Jede Belegung \(\nu\models\operatorname{One}_F(\mathbf V)\)
    induziert das \emph{kodierte Labelling} \(\lambda_\nu\) durch
    \[
        \lambda_\nu(a)=
        \begin{cases}
            \mathsf{in},  &\text{falls }\nu(i_a)=1,\\
            \mathsf{out}, &\text{falls }\nu(o_a)=1,\\
            \mathsf{und}, &\text{falls }\nu(u_a)=1.
        \end{cases}
    \]
    Mit \(\mathbf V^k=(I^k,O^k,U^k)\) bezeichnen wir eine frische
    \emph{\(k\)-te Kopie} dieses Blocks, deren Variablen mit dem
    Superskript \(k\) versehen werden. Ein Quantor über einem Block ist
    die Abkürzung für die Folge gleichartiger Quantoren über allen
    Variablen \(i_a,o_a,u_a\) mit \(a\in A\); insbesondere bindet
    \(\forall\mathbf V\) sämtliche Variablen des Blocks universell.
\end{Definition}

\subsection{Complete- und Stable-Semantik}

\begin{Definition}
    \label{def:co-st-encodings}
    Sei \(F=(A,R)\) ein Argumentationssystem und
    \(\mathbf V=(I,O,U)\) ein Labellingvariablenblock. Die
    \emph{Complete-Kodierung} und die \emph{Stable-Kodierung} sind
    \[
        \operatorname{Co}_F(\mathbf V)
        :=
        \operatorname{One}_F(\mathbf V)
        \land
        \bigwedge_{a\in A}
        \left(
            i_a\leftrightarrow
            \bigwedge_{b\in a^-}o_b
        \right)
        \land
        \bigwedge_{a\in A}
        \left(
            o_a\leftrightarrow
            \bigvee_{b\in a^-}i_b
        \right)
    \]
    beziehungsweise
    \[
        \operatorname{St}_F(\mathbf V)
        :=
        \operatorname{Co}_F(\mathbf V)
        \land
        \bigwedge_{a\in A}\neg u_a.
    \]
\end{Definition}

Die beiden Äquivalenzen der Complete-Kodierung entsprechen unmittelbar den
beiden Bedingungen eines vollständigen Labellings. Die Stable-Kodierung
schließt zusätzlich das Label \(\mathsf{und}\) aus.

\subsection{Preferred-Semantik}

Die Preferred-Semantik ergänzt die Complete-Semantik um eine
Maximalitätsbedingung. Hierfür wird eine zweite Kopie des Variablenblocks
benötigt.

\begin{Definition}
    \label{def:strict-inclusion-and-dominator}
    Seien \(\mathbf V=(I,O,U)\) und
    \(\mathbf V'=(I',O',U')\) zwei Labellingvariablenblöcke für
    \(F=(A,R)\). Die Formel
    \[
        I\prec I'
        :=
        \left(
            \bigwedge_{a\in A}(i_a\rightarrow i'_a)
        \right)
        \land
        \left(
            \bigvee_{a\in A}(\neg i_a\land i'_a)
        \right)
    \]
    heißt \emph{strikte Inklusionsformel}. Sie ist genau dann wahr, wenn
    die durch \(I\) kodierte Argumentmenge eine echte Teilmenge der durch
    \(I'\) kodierten Argumentmenge ist.

    Die Formel
    \[
        \operatorname{Dom}_F(\mathbf V,\mathbf V')
        :=
        \operatorname{Co}_F(\mathbf V')\land(I\prec I')
    \]
    heißt \emph{Dominanzformel}. Eine sie erfüllende Belegung von
    \(\mathbf V'\) heißt \emph{Dominanzzeuge} für \(\mathbf V\).
\end{Definition}

\begin{Definition}
    \label{def:preferred-encoding}
    Sei \(\mathbf V\) ein Labellingvariablenblock und sei
    \(\widehat{\mathbf V}\) ein frischer Block. Die offene QBF
    \[
        \operatorname{Pr}_F(\mathbf V)
        :=
        \operatorname{Co}_F(\mathbf V)
        \land
        \forall\widehat{\mathbf V}\,
        \neg\operatorname{Dom}_F
        (\mathbf V,\widehat{\mathbf V})
    \]
    heißt \emph{Preferred-Kodierung}. Ihre freien Variablen sind die
    Variablen aus \(\mathbf V\).
\end{Definition}

\begin{Definition}
    \label{def:uniform-semantics-predicate}
    Für
    \(\sigma\in\{\mathsf{co},\mathsf{st},\mathsf{pr}\}\) bezeichnet
    die \emph{Semantikkodierung}
    \[
        \operatorname{Sem}^{\sigma}_F(\mathbf V)
        :=
        \begin{cases}
            \operatorname{Co}_F(\mathbf V),
                &\text{falls }\sigma=\mathsf{co},\\
            \operatorname{St}_F(\mathbf V),
                &\text{falls }\sigma=\mathsf{st},\\
            \operatorname{Pr}_F(\mathbf V),
                &\text{falls }\sigma=\mathsf{pr}.
        \end{cases}
    \]
\end{Definition}

\begin{Lemma}
    \label{lem:semantics-encoding-correct}
    Sei \(\nu\models\operatorname{One}_F(\mathbf V)\), und sei
    \(\lambda_\nu\) das durch \(\nu\) kodierte Labelling. Für jedes
    \(\sigma\in\{\mathsf{co},\mathsf{st},\mathsf{pr}\}\) gilt
    \[
        \nu\models\operatorname{Sem}^{\sigma}_F(\mathbf V)
        \quad\Longleftrightarrow\quad
        \lambda_\nu\in\Lambda_\sigma(F).
    \]
\end{Lemma}

\begin{proof}
    Für \(\sigma=\mathsf{co}\) geben die beiden Äquivalenzen in
    \(\operatorname{Co}_F\) genau die Definition eines vollständigen
    Labellings wieder. Für \(\sigma=\mathsf{st}\) kommt die Bedingung
    \(\mathsf{und}(\lambda_\nu)=\emptyset\) hinzu.

    Für \(\sigma=\mathsf{pr}\) ist
    \(\operatorname{Dom}_F(\mathbf V,\widehat{\mathbf V})\) genau dann
    wahr, wenn \(\widehat{\mathbf V}\) ein vollständiges Labelling
    kodiert, dessen In-Menge die In-Menge von \(\lambda_\nu\) echt
    enthält. Die universelle Negation solcher Dominanzzeugen ist daher
    äquivalent zur In-Maximalität von \(\lambda_\nu\) unter allen
    vollständigen Labellings.
\end{proof}

Alle verwendeten Formeln besitzen eine Größe in
\(O(|A|+|R|)\). Insbesondere wird jede Angriffsbeziehung nur in der
Kodierung des jeweils angegriffenen Arguments verwendet.

\section{Einheitliche QBF-Kodierung der Unabhängigkeit}
\label{sec:uniform-independence-qbf}

Die Definition der Unabhängigkeit benötigt drei Labellings. Ihre
Übereinstimmung auf ausgewählten Argumentmengen wird unabhängig von der
Semantik kodiert.

\begin{Definition}
    \label{def:equality-and-mix-formulas}
    Seien \(\mathbf V^j\) und \(\mathbf V^k\) zwei Kopien des
    Labellingvariablenblocks und sei \(S\subseteq A\). Die
    \emph{Gleichheitsformel auf \(S\)} ist
    \[
        \operatorname{Eq}^{j,k}_S
        :=
        \bigwedge_{a\in S}
        \Bigl(
            (i_a^j\leftrightarrow i_a^k)
            \land(o_a^j\leftrightarrow o_a^k)
            \land(u_a^j\leftrightarrow u_a^k)
        \Bigr).
    \]
    Für paarweise disjunkte Mengen \(X,Y,Z\subseteq A\) heißt
    \[
        \operatorname{Mix}^{1,2,3}_{X,Y\mid Z}
        :=
        \operatorname{Eq}^{1,3}_X
        \land
        \operatorname{Eq}^{2,3}_Y
        \land
        \operatorname{Eq}^{1,3}_Z
    \]
    die \emph{Kombinationsformel} der drei Blöcke.
\end{Definition}

\begin{Lemma}
    \label{lem:equality-and-mix-correct}
    Kodieren die Blöcke \(\mathbf V^1,\mathbf V^2,\mathbf V^3\) die
    Labellings \(\lambda_1,\lambda_2,\lambda_3\), so gilt
    \[
        \operatorname{Eq}^{j,k}_S
        \quad\Longleftrightarrow\quad
        \lambda_j|_S=\lambda_k|_S
    \]
    und
    \[
        \operatorname{Mix}^{1,2,3}_{X,Y\mid Z}
        \quad\Longleftrightarrow\quad
        \left(
        \lambda_3|_X=\lambda_1|_X
        \land
        \lambda_3|_Y=\lambda_2|_Y
        \land
        \lambda_3|_Z=\lambda_1|_Z
        \right).
    \]
\end{Lemma}

\begin{proof}
    Aufgrund der Wohlgeformtheit ist für jedes Argument genau eine der drei
    Labelvariablen wahr. Die drei Äquivalenzen für ein Argument sind daher
    genau dann erfüllt, wenn beide Blöcke dasselbe Label kodieren. Die zweite
    Aussage folgt durch Einsetzen der drei Gleichheitsformeln.
\end{proof}

\begin{Definition}
    \label{def:generic-independence-qbf}
    Sei \(F=(A,R)\), seien \(X,Y,Z\subseteq A\) paarweise disjunkt und sei
    \(\sigma\in\{\mathsf{co},\mathsf{st},\mathsf{pr}\}\). Die
    \emph{semantikparametrisierte Unabhängigkeits-QBF} ist
    \[
    \begin{split}
        \Phi^\sigma_F(X,Y\mid Z)
        :=\forall\mathbf V^1\,\forall\mathbf V^2\,
        \Bigl(&
            \operatorname{Sem}^{\sigma}_F(\mathbf V^1)
            \land
            \operatorname{Sem}^{\sigma}_F(\mathbf V^2)
            \land
            \operatorname{Eq}^{1,2}_Z
        \\
        &\rightarrow
        \exists\mathbf V^3\,
        \bigl(
            \operatorname{Sem}^{\sigma}_F(\mathbf V^3)
            \land
            \operatorname{Mix}^{1,2,3}_{X,Y\mid Z}
        \bigr)
        \Bigr).
    \end{split}
    \]
\end{Definition}

\begin{Theorem}
    \label{thm:generic-independence-qbf-correct}
    Für jedes \(F=(A,R)\), alle paarweise disjunkten Mengen
    \(X,Y,Z\subseteq A\) und jedes
    \(\sigma\in\{\mathsf{co},\mathsf{st},\mathsf{pr}\}\) gilt
    \[
        \Phi^\sigma_F(X,Y\mid Z)\text{ ist wahr}
        \quad\Longleftrightarrow\quad
        X\perp_\sigma Y\mid Z.
    \]
\end{Theorem}

\begin{proof}
    Sei die QBF wahr, und seien
    \(\lambda_1,\lambda_2\in\Lambda_\sigma(F)\) mit
    \(\lambda_1|_Z=\lambda_2|_Z\) gegeben. Ihre kodierenden Belegungen
    erfüllen nach Lemma~\ref{lem:semantics-encoding-correct} die beiden
    Semantikkodierungen und nach
    Lemma~\ref{lem:equality-and-mix-correct} die Formel
    \(\operatorname{Eq}^{1,2}_Z\). Die QBF liefert daher eine Belegung von
    \(\mathbf V^3\), die ein \(\sigma\)-Labelling \(\lambda_3\) kodiert
    und die geforderten Restriktionen erfüllt. Somit gilt
    \(X\perp_\sigma Y\mid Z\).

    Gelte umgekehrt \(X\perp_\sigma Y\mid Z\), und seien die beiden
    universellen Blöcke beliebig belegt. Ist das Antezedens falsch, ist die
    Implikation wahr. Andernfalls kodieren die Blöcke zwei
    \(\sigma\)-Labellings, die auf \(Z\) übereinstimmen. Nach Definition
    der Unabhängigkeit existiert ein passendes \(\lambda_3\). Seine
    kodierende Belegung erfüllt das Konsequens. Dies gilt für jede Belegung
    der universellen Blöcke, also ist die QBF wahr.
\end{proof}

\subsection{Quantorenstruktur der Spezialfälle}
\label{subsec:independence-prenex}

Für Complete und Stable sind die Semantikkodierungen quantorenfrei. Daher
kann der Existenzquantor für den dritten Block vor die Matrix gezogen
werden.

\begin{Proposition}
    \label{prop:co-st-independence-prenex}
    Für \(\sigma\in\{\mathsf{co},\mathsf{st}\}\) ist
    \(\Phi^\sigma_F(X,Y\mid Z)\) logisch äquivalent zu
    \[
    \begin{split}
        \forall\mathbf V^1\,\forall\mathbf V^2\,
        \exists\mathbf V^3\,
        \Bigl(&
        \bigl(
            \operatorname{Sem}^{\sigma}_F(\mathbf V^1)
            \land
            \operatorname{Sem}^{\sigma}_F(\mathbf V^2)
            \land
            \operatorname{Eq}^{1,2}_Z
        \bigr)
        \\
        &\rightarrow
        \bigl(
            \operatorname{Sem}^{\sigma}_F(\mathbf V^3)
            \land
            \operatorname{Mix}^{1,2,3}_{X,Y\mid Z}
        \bigr)
        \Bigr).
    \end{split}
    \]
    Die Formel besitzt damit zwei alternierende Quantorenblöcke.
\end{Proposition}

Bei Preferred dürfen die internen Quantoren der drei
Maximalitätsprüfungen nicht lediglich nebeneinandergestellt werden. Die
Maximalität der universell gewählten Labellings steht im Antezedens und
kehrt deshalb beim Übergang in Negationsnormalform ihre Quantorenrichtung
um. Die folgende Formel macht diese Abhängigkeit explizit.

\begin{Definition}
    \label{def:preferred-independence-prenex}
    Für \(k\in\{1,2,3\}\) sei
    \(\widehat{\mathbf V}^{k}\) ein frischer Labellingvariablenblock. Die
    \emph{pränexe Preferred-Unabhängigkeitskodierung} ist
    \[
    \begin{split}
        \Phi^{\mathsf{pr},\mathrm{PNF}}_F(X,Y\mid Z)
        := {}&
        \forall\mathbf V^1\,\forall\mathbf V^2\,
        \exists\mathbf V^3\,
        \exists\widehat{\mathbf V}^{1}\,
        \exists\widehat{\mathbf V}^{2}\,
        \forall\widehat{\mathbf V}^{3}
        \;\Bigl[
        \\
        &\neg\operatorname{Co}_F(\mathbf V^1)
        \lor
        \neg\operatorname{Co}_F(\mathbf V^2)
        \lor
        \neg\operatorname{Eq}^{1,2}_Z
        \\
        &\lor
        \operatorname{Dom}_F
            (\mathbf V^1,\widehat{\mathbf V}^{1})
        \lor
        \operatorname{Dom}_F
            (\mathbf V^2,\widehat{\mathbf V}^{2})
        \\
        &\lor
        \Bigl(
            \operatorname{Co}_F(\mathbf V^3)
            \land
            \operatorname{Mix}^{1,2,3}_{X,Y\mid Z}
            \land
            \neg\operatorname{Dom}_F
                (\mathbf V^3,\widehat{\mathbf V}^{3})
        \Bigr)
        \Bigr].
    \end{split}
    \]
\end{Definition}

\begin{Proposition}
    \label{prop:preferred-independence-prenex-correct}
    Es gilt
    \[
        \Phi^{\mathsf{pr},\mathrm{PNF}}_F(X,Y\mid Z)
        \quad\Longleftrightarrow\quad
        \Phi^{\mathsf{pr}}_F(X,Y\mid Z).
    \]
    Insbesondere ist die pränexe Formel genau dann wahr, wenn
    \(X\perp_{\mathsf{pr}}Y\mid Z\) gilt.
\end{Proposition}

\begin{proof}
    Fixiere Belegungen von \(\mathbf V^1\) und \(\mathbf V^2\). Kodiert
    einer der Blöcke kein vollständiges Labelling oder stimmen die Blöcke
    auf \(Z\) nicht überein, ist eine der ersten drei Disjunktionen wahr.
    Kodiert einer der Blöcke ein vollständiges, aber nicht bevorzugtes
    Labelling, kann der Existenzspieler in
    \(\widehat{\mathbf V}^{1}\) beziehungsweise
    \(\widehat{\mathbf V}^{2}\) einen Dominanzzeugen wählen.

    Sind dagegen beide Blöcke bevorzugt und stimmen auf \(Z\) überein,
    bleiben alle genannten Disjunktionen falsch. Dann muss ein Block
    \(\mathbf V^3\) gewählt werden, der vollständig ist und die
    Kombinationsformel erfüllt. Die anschließende Allquantifizierung über
    \(\widehat{\mathbf V}^{3}\) verlangt zusätzlich, dass kein
    Dominanzzeuge für \(\mathbf V^3\) existiert. Dies ist genau die
    Preferred-Bedingung im Konsequens der generischen Formel.
\end{proof}

\section{Komplexität}
\label{sec:independence-complexity}

Die Quantorenstruktur der Kodierungen entspricht der bekannten
Komplexitätsklassifikation. Zum Vergleich wird zunächst die Komplexität der
Prüfung eines einzelnen Labellings festgehalten.

\begin{Definition}
    \label{def:labelling-verification-problem}
    Für eine Semantik \(\sigma\) bezeichnet
    \(\mathrm{VER}_\sigma\) das \emph{Labelling-Verifikationsproblem}:
    Eine Instanz besteht aus einem Argumentationssystem \(F=(A,R)\) und
    einem Labelling \(\lambda\colon A\to
    \{\mathsf{in},\mathsf{out},\mathsf{und}\}\). Gefragt ist, ob
    \(\lambda\in\Lambda_\sigma(F)\) gilt.
\end{Definition}

Nach der von \textcite[Tabelle~3.1]{dvorak:2012:computational-aspects}
zusammengefassten Klassifikation liegen
\(\mathrm{VER}_{\mathsf{co}}\) und
\(\mathrm{VER}_{\mathsf{st}}\) in \(\mathrm P\), während
\(\mathrm{VER}_{\mathsf{pr}}\) \(\mathrm{coNP}\)-vollständig ist. In
der QBF-Kodierung zeigt sich dieser Unterschied in der zusätzlichen
Maximalitätsquantifizierung der Preferred-Semantik.

\textcite[Proposition~C.4]{bluemel:2025:independence-aba} zeigen das
folgende Resultat für abstrakte Argumentationssysteme.

\begin{Theorem}
    \label{thm:independence-complexity}
    Bezüglich polynomialer Many-one-Reduktionen sind
    \[
        \mathrm{IND}_{\mathsf{co}}
        \quad\text{und}\quad
        \mathrm{IND}_{\mathsf{st}}
    \]
    \(\Pi^{\mathrm P}_2\)-vollständig. Das Problem
    \(\mathrm{IND}_{\mathsf{pr}}\) ist
    \(\Pi^{\mathrm P}_3\)-vollständig.
\end{Theorem}

\begin{proof}
    Für Complete und Stable liefert
    Proposition~\ref{prop:co-st-independence-prenex} in Polynomialzeit eine
    äquivalente QBF mit Präfix \(\forall\exists\) und polynomiell großer
    Matrix. Daher liegen beide Probleme in \(\Pi^{\mathrm P}_2\). Für
    Preferred besitzt die äquivalente Formel aus
    Definition~\ref{def:preferred-independence-prenex} das Präfix
    \(\forall\exists\forall\) und eine polynomiell große Matrix; folglich
    liegt das Problem in \(\Pi^{\mathrm P}_3\).

    Für die Härte reduzieren
    \textcite[Propositionen~C.1--C.4]
    {bluemel:2025:independence-aba} polynomial von
    \(\mathrm{QBF}^{\forall}_{2,\mathrm{KNF}}\) auf die Complete- und
    Stable-Probleme. Für Preferred wird polynomial vom Falschheitsproblem
    für Formeln der Gestalt
    \(\exists X\,\forall Y\,\exists Z\,\varphi\) mit KNF-Matrix, also von
    \(\overline{\mathrm{QBF}^{\exists}_{3,\mathrm{KNF}}}\), reduziert.
    Die verwendeten Ausgangsprobleme sind die kanonischen
    \(\Pi^{\mathrm P}_2\)- beziehungsweise
    \(\Pi^{\mathrm P}_3\)-vollständigen QBF-Probleme. Damit folgen die
    behaupteten Härten und zusammen mit den Zugehörigkeiten die
    Vollständigkeiten.
\end{proof}

\begin{table}[tb]
    \centering
    \caption{Einheitlicher Vergleich von Verifikation, QBF-Präfix und
    Unabhängigkeitskomplexität. Die Verifikationskomplexität folgt
    \textcite[Tabelle~3.1]{dvorak:2012:computational-aspects}, die
    Vollständigkeit der Unabhängigkeitsprobleme
    \textcite[Proposition~C.4]{bluemel:2025:independence-aba}.}
    \label{tab:independence-complexity-comparison}
    \small
    \begin{tabular}{llll}
        \toprule
        Semantik
        & \(\mathrm{VER}_\sigma\)
        & Präfix der Unabhängigkeits-QBF
        & \(\mathrm{IND}_\sigma\) \\
        \midrule
        \(\mathsf{co}\)
        & in \(\mathrm P\)
        & \(\forall\mathbf V^1,\mathbf V^2\,
             \exists\mathbf V^3\)
        & \(\Pi^{\mathrm P}_2\)-vollständig \\
        \(\mathsf{st}\)
        & in \(\mathrm P\)
        & \(\forall\mathbf V^1,\mathbf V^2\,
             \exists\mathbf V^3\)
        & \(\Pi^{\mathrm P}_2\)-vollständig \\
        \(\mathsf{pr}\)
        & \(\mathrm{coNP}\)-vollständig
        & \(\forall\mathbf V^1,\mathbf V^2\,
             \exists\mathbf V^3,\widehat{\mathbf V}^{1},
             \widehat{\mathbf V}^{2}\,
             \forall\widehat{\mathbf V}^{3}\)
        & \(\Pi^{\mathrm P}_3\)-vollständig \\
        \bottomrule
    \end{tabular}
\end{table}

\section{Empirischer Vergleich der Semantiken}
\label{sec:empirical-independence-occurrence}